\chapter{Cosmologie de l'état généalogique global : connexion, courbure, expansion et mémoire de frontière}
\label{chap:cosmologia}

\section{La totalité comme état compatible}

Une cosmologie HMT--MD est une réalisation globale du système de chemins,
de frontières et de mémoire. Son état appartient à la limite des sections compatibles ;
sa géométrie s'exprime au moyen d'une cotétrade, d'une connexion et d'une métrique ; son évolution
publie facteur d'échelle, courbure, contenu matériel et horizon. La
totalité intervient comme objet mathématique car chaque section locale est une
restriction d'un état global avec des conditions de clôture.

Cette idée offre une formulation précise du principe de Mach. Soient
\(b:\mathcal X_\infty\to B_\infty\) la trace globale de frontière et
\(I_v:\mathcal X_\infty\to V_v\) une réponse inertielle dans un voisinage
\(v\). La réponse est machienne lorsqu'il existe
\[
 \overline I_v:\operatorname{im}(b)\to V_v
\]
tel que
\[
 I_v=\overline I_v\circ b,
\]
et lorsque la restriction purement locale laisse ouvertes des réponses distinctes.
L'application \(\overline I_v\) existe de manière unique exactement lorsque
\(I_v\) est constante sur chaque fibre de \(b\). L'inertie locale est ainsi
déterminée par la classe globale de clôture.

Le référentiel d'une section s'obtient en restreignant simultanément l'état et
la connexion :
\[
 \mathfrak F_v(x)
 =\operatorname{Res}_v(x,\nabla_{\rm TPK}).
\]
Une géodésique affine satisfait
\[
 \nabla_{\dot\gamma}\dot\gamma=0,
\]
et un référentiel parallèle satisfait
\(\nabla_{\dot\gamma}e_a=0\). Accélération, rotation, torsion, courbure et
holonomie mesurent les défauts par rapport à ce transport induit.

\section{\(3\) cartes de courbure}

Les régimes elliptique, parabolique et hyperbolique du TRIT fournissent trois
domaines algébriques. Leur réalisation géométrique exige une application
\[
 \mathcal M_{s,\rho}:\mathcal A_s
 \longrightarrow(\mathcal X_{s,\rho},g_{s,\rho})
\]
qui préserve connexion, courbure, orientation et données de bord. Cette
interface étant déclarée, l'atlas coordonné utilise une carte de courbure
négative de type anti--de Sitter, une carte euclidienne gnomonique et une carte de
courbure positive de type de Sitter.

La carte négative organise le carré holographique entre état intérieur et
observables conformes de frontière. La carte euclidienne réalise la descente locale
\[
 \frac{\mathcal G_L}{\mathcal I_L}=1
\]
du principe d'équivalence. La carte positive reçoit l'expansion
accélérée lorsque sa dynamique satisfait
\[
 \frac{\ddot a}{a}=H^2>0.
\]
L'antécédent commun conserve le rapport aréal--volumétrique
\[
 \frac{\mathcal L_{\rm ar}}{\mathcal L_{\rm vol}}
 =\frac{\pi_{\rm HMT}}{\varphi_{\rm HMT}^2}.
\]
La correspondance entre signe algébrique et géométrie physique se démontre dans
chaque carte au moyen de l'application métrique, de sorte que l'atlas conserve le type de
chaque transition.

La géométrie de de Sitter fournit un résultat exact de comparaison.
Pour \(H>0\), ses feuilletages fermé, plat et ouvert possèdent
\[
 a_{+1}(t)=H^{-1}\cosh(Ht),\qquad
 a_0(t)=a_*e^{Ht},\qquad
 a_{-1}(t)=H^{-1}\sinh(Ht),
\]
dans leurs domaines respectifs. Les trois satisfont
\[
 \left(\frac{\dot a_k}{a_k}\right)^2+\frac{k}{a_k^2}=H^2,
 \qquad
 \frac{\ddot a_k}{a_k}=H^2,
\]
\[
 R^{(4)}=12H^2,\qquad
 R^{(3)}=\frac{6k}{a_k^2}.
\]
Le signe \(k\) caractérise la courbure de la tranche spatiale ; la
courbure de l'espace-temps complet reste positive et commune.

\section{Mémoire distribuée et expansion}

Le transport dodécaphasique sépare une partie paire et une partie impaire. La partie
paire conserve le coût partagé par un chemin et son inverse ; l'impaire conserve
l'orientation. Dans la lecture cosmologique, le secteur pair peut se réaliser
comme réponse moyenne à l'expansion et le secteur impair comme orientation du
feuillet.

Une mémoire distribuée peut provenir d'un générateur local. Le noyau
\[
 K_q(m,n)=q^{|m-n|},\qquad 0<q<1,
\]
est positif et son inverse est tridiagonale. Les corrélations entre profondeurs
arbitraires sont donc le Green d'une action de premiers voisins. Cette
identité permet à l'histoire globale de conserver la mémoire tandis que la loi de
mise à jour agit localement.

Dans une géométrie homogène, la réponse paire admet la forme constitutive
\[
 p_{\rm eff}=p-3\zeta H.
\]
Le coefficient \(\zeta\) représente la résistance de la mémoire symétrique à
l'expansion. Sa sélection appartient à la réalisation cosmologique et doit
descendre de l'opérateur de transport et de sa carte dimensionnelle.

\section{Complétion (1000=729+243+27+1), données de frontière et critère de rebond d'Einstein–Cartan}

La complétion
\[
 1000=729+243+27+1
\]
organise intérieur, strates de frontière et sceau. Dans une cosmologie, ce
motif se réalise sur une paire \((M,\partial M)\) avec cotétrade, connexion,
matière et données de bord. La frontière reçoit un caractère spatial, temporel ou
nul ; l'action incorpore ses termes de bord et de coin ; les conditions de
raccordement fixent les histoires qui se prolongent.

Un rebond d'Einstein--Cartan peut apparaître lorsque la matière à spin
produit une contribution effective qui croît comme
\[
 s^2\propto a^{-6}.
\]
Dans une convention de densité de masse, l'équation prend le schéma
\[
 H^2=\frac{8\pi G}{3}
 \bigl(\rho_{\rm m}-\rho_{\rm spin,m}\bigr)+\cdots.
\]
Le point d'équilibre
\(\rho_{\rm m}=\rho_{\rm spin,m}\) est candidat à un minimum régulier du
facteur d'échelle lorsque l'action de matière, le signe, les conditions
d'énergie et le problème de frontière le soutiennent. HMT ajoute l'obligation de
dériver le courant de spin et le seuil à partir de l'état enrichi.

Une transition vers des composantes cosmologiques descendantes est représentée
au moyen d'un cobordisme
\[
 \mathcal W:
 \Sigma_{\rm in}\longrightarrow
 \Sigma_{\rm prog}\sqcup\Sigma_{\rm desc}.
\]
La conservation de l'information s'exprime au moyen d'une charge cocyclique, d'une
isométrie ou d'une évolution unitaire. Dans le cas quantique, les entropies,
l'information mutuelle et les corrélations remplacent une somme scalaire simple.
Les termes de raccordement, la causalité, l'hyperbolicité et le contrôle des
singularités font partie de l'objet.

\section{Système commun d'observables cosmologiques : expansion, lentilles, croissance structurelle et anisotropies du fond}

Une réalisation destinée à expliquer les divergences attribuées à la matière
noire ou à l'énergie noire doit utiliser un ensemble commun de paramètres pour
les courbes de rotation, les lentilles gravitationnelles, la croissance des structures,
les anisotropies du fond cosmique et l'histoire de l'expansion. Cette exigence
transforme la cosmologie des chemins en un système conjoint d'observables
corrélés soumis à une confrontation simultanée.

Le cycle dodécaphasique apporte en outre un sélecteur harmonique. Pour
\[
 w_k=\frac1{12}
 \left[1+2\varepsilon\cos(3\phi_k-\delta)\right],
\]
la transformée discrète possède un support uniquement dans
\(m=0,\pm3\pmod{12}\). Une prédiction céleste exige de fixer avant la confrontation
l'application vers le ciel, l'axe, l'amplitude, la phase, le masque et la statistique.
La comparaison en aveugle conserve la distinction entre le théorème harmonique et sa
réalisation cosmologique.

\begin{resultado}
La factorisation machienne, le critère d'auto-inertie, les trois feuilletages
de de Sitter, la mémoire locale avec Green distribué et le support harmonique
dodécaphasique sont des résultats exacts dans leurs domaines. L'atlas métrique,
l'équation effective d'échelle, le rebond et le cobordisme constituent des interfaces
physiques typées avec des conditions explicites et composent la cosmologie des chemins.
\end{resultado}

\begin{lecturafisica}
L'univers est décrit comme une section globale dont la frontière conditionne les
référentiels locaux. L'expansion transforme la mémoire collective ; la courbure et
les ondes renvoient cette transformation aux chemins ; l'horizon enregistre la
partie accessible de l'histoire.
\end{lecturafisica}

\begin{frontera}
La réalisation serait réfutée si la réponse locale variait dans une
même fibre globale, si l'application métrique ne respectait pas la courbure ou l'orientation,
si l'accélération était dépourvue d'équation dynamique, si le rebond produisait une
singularité ou violait les conditions de raccordement, ou si un unique ensemble de
paramètres était incompatible avec expansion, lentilles, structure et fond
cosmique.
\end{frontera}
